% Template-independent section text; included from paper/main.tex once the venue template is in place.
\section{Setting and method}
\label{sec:method}

\paragraph{Leaderboard and aggregation rules.}
A leaderboard is a score matrix $X\in\mathbb{R}^{M\times n}$ (higher is better) of $M$ models on $n$ sub-datasets (``tasks''),
and an aggregation rule $R$ that maps $X$ to a ranking of the models.
Rules are either \emph{score rules}, which compute a scalar $f(x_j)$ from each model's score vector
(mean, geometric, harmonic and power means with $p<1$, median, minimum, trimmed mean),
or \emph{rank rules}, which depend only on per-task orderings
(Borda/mean rank, plurality as first-place count, and the requirement of being a Condorcet winner).
Being a Condorcet winner is sufficient for first place under every Condorcet-consistent rule
(Copeland, Schulze, Kemeny, ranked pairs), so its minimal investment is an upper bound on theirs.

\paragraph{Targeted investment.}
A developer wants a \emph{target} model $i$ to rank first.
Training on data related to task $t$ lets them raise the target's score on $t$ from $x_{it}$ to any $y_t\in[x_{it},c_t]$,
where $c_t$ is a cap (headroom); all other models are held fixed.
The target is \emph{first} if it is strictly above every other model under $R$ (margin $\varepsilon=10^{-3}$ score points).
A cap scenario is the central modelling assumption: we study $c_t=100$ (ceil), the best existing score on $t$ plus $0$ or $2$ points
(front), a fixed gain of $5/10/20$ points over the target's own score, and empirically anchored caps
equal to the 95th/99th percentile of fine-tune-minus-base differences measured on the Open LLM Leaderboard.
\emph{Minimal investment} is $\min_y \mathrm{cost}(y-x_i)$ over feasible $y$ making $i$ first.
We report several costs side by side:
$K$, the number of tasks touched (breadth); $L_1=\sum_t (y_t-x_{it})$ (volume);
$L_\infty=\max_t (y_t-x_{it})$; $L_2=\sum_t (y_t-x_{it})^2$ (a convex cost under which concentrating on one task is no longer free);
and $u^\ast$, the smallest uniform gain $\delta$ on every task that makes $i$ first, the ``general improvement'' reference.
If no feasible $y$ exists the target is \emph{infeasible}; this is an outcome that we report as a fraction and never drop or impute.

\paragraph{Exact solvers.}
For each rule we compute minimal investment exactly rather than by heuristic search.
For the mean, $L_1^\ast=n(\tau+\varepsilon)-\sum_t x_{it}$ in closed form, where $\tau$ is the best rival's value, and $K^\ast$ follows by raising the tasks with the most headroom first.
For geometric, harmonic and power means ($p<1$, concave) the $L_1$ optimum equalises the shifted scores at a water level, found by bisection; $K$ is a greedy selection.
For median and min the optimum is an order-statistic argument (how many tasks must reach the rivals' decisive score).
The trimmed mean and the rank rules are solved as mixed-integer programs with HiGHS: a raise on task $t$ to just above a rival's score passes that rival on $t$ (a tie counts one half), so the candidate levels are finite, and rival constraints are generated lazily.
Being the Condorcet winner is the requirement that the target beat every rival in a majority of tasks; since it is sufficient for first place under every Condorcet-consistent rule, we also solve the Copeland rule exactly on HELM to see how tight this bound is.
Every solver is tested against brute-force enumeration on small random instances, and the closed forms for the mean, median and min are compared against the solver output on every target (Section~\ref{sec:exp-theory}).
A solve that reaches its 120\,s limit is flagged and never recorded as infeasible; none did in any file that records the flag (Appendix~\ref{app:repro}).
